\begin{figure}[ht]
\centering
\begin{tikzpicture}[x=1cm, y=1cm, font=\small]

% Mildly elongated quadratic f(x,y) = 1/2 (a x^2 + b y^2) with a=1, b=6.
% Level set f = c is the ellipse x^2/(2c) + y^2/(2c/6) = 1,
% semi-axes (sqrt(2c), sqrt(c/3)). Centre at (3, 2) so axes are visible.
\def\cx{3}
\def\cy{2}

% Axes
\draw[->, thick] (-0.2, 0) -- (6.4, 0) node[right, font=\scriptsize] {$x_{1}$};
\draw[->, thick] (0, -0.2) -- (0, 4.2) node[above, font=\scriptsize] {$x_{2}$};

% Elliptic contours centred at the optimum (cx, cy)
\foreach \c in {0.3, 0.9, 1.8, 3.0, 4.5} {
  \pgfmathsetmacro{\axA}{sqrt(2*\c)}
  \pgfmathsetmacro{\axB}{sqrt(\c/3)}
  \draw[engblue, thin] (\cx, \cy) ellipse ({\axA} and {\axB});
}

% Optimum
\fill[engred] (\cx, \cy) circle (2pt);
\node[engred, anchor=west, font=\scriptsize] at (\cx+0.12, \cy) {$\vect{x}^{*}$};

% Gradient-descent iterates with a modest step size: smooth descent
% from x_0 = (0.4, 3.6) curving into the optimum, with the
% characteristic perpendicular-to-contour first steps.
\fill[engblack] (0.40, 3.60) circle (1.5pt);
\fill[engblack] (1.30, 2.55) circle (1.5pt);
\fill[engblack] (1.95, 2.25) circle (1.5pt);
\fill[engblack] (2.45, 2.10) circle (1.5pt);
\fill[engblack] (2.78, 2.04) circle (1.5pt);
\fill[engblack] (2.95, 2.01) circle (1.5pt);

\draw[engblack, thick]
  (0.40, 3.60) -- (1.30, 2.55) -- (1.95, 2.25) -- (2.45, 2.10)
  -- (2.78, 2.04) -- (2.95, 2.01);

\node[engblack, anchor=south east, font=\scriptsize] at (0.40, 3.62) {$\vect{x}_{0}$};

% Gradient direction at the first point (negative gradient = descent step)
\draw[->, engred, thick] (0.40, 3.60) -- (0.85, 3.07);
\node[engred, anchor=west, font=\scriptsize] at (0.88, 3.20)
  {$-\grad f(\vect{x}_{0})$};

% Caption note
\node[anchor=north, font=\scriptsize, align=center] at (3.1, -0.5)
  {$f(\vect{x}) = \tfrac{1}{2}\norm{A\vect{x} - \vect{b}}_{2}^{2}$,
   $\kappa = 6$};

\end{tikzpicture}
\caption[Gradient descent on a least-squares contour]{Gradient
descent on the convex quadratic of a least-squares objective,
plotted on its elliptic level sets. Each step moves along the
negative gradient, which at every point is orthogonal to the local
level set. With a well-chosen step size on a mildly conditioned
problem ($\kappa = 6$), the iterates curve steadily into the
optimum $\vect{x}^{*}$ without the zigzag of the badly conditioned
case in \cref{fig:vol02:ch15:gradient-zigzag}. The optimum is the
point where the residual $A\vect{x}^{*} - \vect{b}$ is orthogonal
to the column space of $A$, the geometric reading of the
stationarity condition $A^{\transpose}(A\vect{x}^{*} - \vect{b}) =
\vect{0}$.}
\label{fig:vol02:ch15:contour-descent}
\end{figure}
